% !TeX root = main.tex
\section{Trim-Like Operating Condition and Evaluation Protocol}
\subsection{Trim Criteria}
We define a trim-like operating condition as one where the cycle-averaged streamwise force is near zero and the cycle-averaged lift is near the vehicle weight:
\begin{equation}
\bar{F}_x \approx 0,\qquad \bar{F}_z \approx mg.
\end{equation}
In this paper, we treat trim identification as a model-internal procedure used to establish a reproducible baseline for comparing duty-cycle modulation.

\subsection{Cycle Averaging}
To avoid phase truncation bias, all reported means are computed over an integer number of cycles.
With simulation timestep $\Delta t = 1/240~\mathrm{s}$ and $f=3.8~\mathrm{Hz}$, we use $N=1200$ steps ($5~\mathrm{s}$), corresponding to $19$ cycles.
Cycle-averaged quantities are computed by time integration over the window.

\subsection{Baseline Operating Point}
We identify a baseline operating point by scanning frequency, body pitch, and airspeed while keeping other parameters fixed, and selecting a point that yields $\bar{F}_x$ near zero and lift on the order of $mg$ for the target mass scale.
The baseline used for the duty-cycle sweep in this paper is:
\begin{equation}
f = 3.8~\mathrm{Hz},\quad U = 7.0~\mathrm{m/s},\quad \theta_{\mathrm{body}} = 13^\circ,
\end{equation}
with separation disabled and prescribed twist enabled.
